\documentclass[12pt]{article}

% Packages
\usepackage{amsmath}
\usepackage{amssymb}
\usepackage{amsthm}
\usepackage{geometry}
\usepackage{mathrsfs}
\usepackage{mathtools}
\usepackage{tikz}
\usepackage[colorlinks=true, linkcolor=blue, citecolor=blue, urlcolor=blue]{hyperref}

% Page Geometry
\geometry{letterpaper, margin=1in}

% Theorem Environments
\theoremstyle{definition}
\newtheorem{problem}{Problem}
\renewcommand{\qedsymbol}{$\blacksquare$}

% Title Information
\title{\textbf{Rigorous Practice Problem Set:\\ Fundamental Properties of Set Operations}}
\author{Advanced Mathematics, Logic, \& Quantitative Modeling}
\date{October 2026}

\begin{document}

\maketitle

\tableofcontents
\newpage

% ==========================================
% SECTION 1: CONCEPT CHECKLIST & MAPPING
% ==========================================
\section{Concept Checklist \& Question Mapping}
\label{sec:checklist_start}

This problem set is strictly derived from the theoretical frameworks, logical isomorphisms, and real-world modeling paradigms established in the source text. Every problem maps to one or more of the following core concepts:

\begin{itemize}
    \item \textbf{C1: Element-Chasing Proofs (Mutual Inclusion):} Rigorous deduction of $X \subseteq Y$ and $Y \subseteq X$. (Problems 1, 2, 8, 9)
    \item \textbf{C2: Biconditional Logical Equivalences:} Chaining $\iff$ statements using Boolean logic. (Problems 3, 10)
    \item \textbf{C3: Algebraic Set Manipulation:} Utilizing established axioms (De Morgan's, Distributive) as algebraic variables. (Problems 4, 5, 6, 7)
    \item \textbf{C4: Set Difference Translations:} Converting $A - B$ to $A \cap \overline{B}$. (Problems 4, 5, 6, 11, 12)
    \item \textbf{C5: Propositional Logic to Set Theory Isomorphism:} Translating tautologies and logical laws to set identities. (Problems 13, 14, 15)
    \item \textbf{C6: Infinite Families of Sets:} Arbitrary unions ($\bigcup_{i \in I}$) and intersections ($\bigcap_{i \in I}$). (Problems 16, 17, 18, 19)
    \item \textbf{C7: Generalized De Morgan's Laws:} Extending finite complements to arbitrary index sets. (Problem 16)
    \item \textbf{C8: Probability Spaces \& Measure Theory ($\sigma$-algebras):} Set-theoretic foundations for mutually exclusive events and probability bounds. (Problems 20, 21, 22, 23)
    \item \textbf{C9: Quantitative Finance Modeling:} Binomial pricing, risk-neutral measures ($\mathbb{Q}$), and expected discounted payoffs. (Problems 24, 25, 26)
    \item \textbf{C10: Machine Learning Modeling:} Feature space partitioning and decision tree logic. (Problems 27, 28)
    \item \textbf{C11: Diagnostic Flaws \& Counterexamples:} Identifying inverse fallacies, non-associativity, and vacuous truth errors. (Problems 29, 30)
\end{itemize}

\newpage

% ==========================================
% SECTION 2: PROBLEM SET
% ==========================================
\section{Problem Set}
\label{sec:problems}

\subsection{Fundamental Identities and Element-Chasing}
\begin{problem}
Using the method of element-chasing (mutual inclusion), prove the Absorption Law: $A \cup (A \cap B) = A$.
\end{problem}

\begin{problem}
Using the method of element-chasing, prove that if $A \subseteq B$, then $A \cap B = A$.
\end{problem}

\begin{problem}
Using a chain of biconditional logical equivalences, prove that $(A - B) - C = A - (B \cup C)$.
\end{problem}

\subsection{Algebraic Set Manipulation}
\begin{problem}
Using algebraic set manipulation and the definition of set difference, prove that $A - (A - B) = A \cap B$.
\end{problem}

\begin{problem}
Prove that $A \cap (B - C) = (A \cap B) - (A \cap C)$ using algebraic set manipulation.
\end{problem}

\begin{problem}
Define the symmetric difference as $A \Delta B = (A - B) \cup (B - A)$. Prove algebraically that $A \Delta B = (A \cup B) - (A \cap B)$.
\end{problem}

\begin{problem}
Prove algebraically that $A \cap (B \Delta C) = (A \cap B) \Delta (A \cap C)$.
\end{problem}

\subsection{Logical Isomorphisms and Implications}
\begin{problem}
Establish the logical equivalence of the contrapositive in set theory: Prove via element-chasing that $A \subseteq B \iff \overline{B} \subseteq \overline{A}$.
\end{problem}

\begin{problem}
Translate the logical tautology $P \implies (P \lor Q)$ into its corresponding set-theoretic statement and prove it via element-chasing.
\end{problem}

\begin{problem}
Using biconditional logical equivalences, prove that $\overline{A \cup (B \cap C)} = (\overline{A} \cap \overline{B}) \cup (\overline{A} \cap \overline{C})$.
\end{problem}

\subsection{Counterexamples and Diagnostic Flaws}
\begin{problem}
Find a counterexample to disprove the claim that set difference is associative: $A - (B - C) = (A - B) - C$.
\end{problem}

\begin{problem}
Find a counterexample to disprove the claim: $(A \cup B) - C = A \cup (B - C)$.
\end{problem}

\begin{problem}
Identify the logical flaw in the following statement and provide the correct negation: ``If $x \notin A \cap B$, then $x \notin A \land x \notin B$.''
\end{problem}

\begin{problem}
Identify the logical flaw in the following proof: \\
\textit{Claim:} If $A \subseteq B \cup C$, then $A \subseteq B$ or $A \subseteq C$. \\
\textit{Proof:} Let $x \in A$. Since $A \subseteq B \cup C$, $x \in B \cup C$. This means $x \in B$ or $x \in C$. Therefore, all elements of $A$ are in $B$, or all elements of $A$ are in $C$. Thus $A \subseteq B$ or $A \subseteq C$.
\end{problem}

\begin{problem}
Prove or disprove: For any sets $A$ and $B$, the power set satisfies $\mathcal{P}(A \cup B) = \mathcal{P}(A) \cup \mathcal{P}(B)$.
\end{problem}

\subsection{Infinite Families of Sets}
\begin{problem}
Let $I$ be an arbitrary index set and $\{A_i\}_{i \in I}$ be a family of sets. Prove the Generalized De Morgan's Law: $\overline{\bigcap_{i \in I} A_i} = \bigcup_{i \in I} \overline{A_i}$.
\end{problem}

\begin{problem}
Prove the Generalized Distributive Law for an arbitrary family of sets $\{A_i\}_{i \in I}$ and a set $B$: $B \cap \left(\bigcup_{i \in I} A_i\right) = \bigcup_{i \in I} (B \cap A_i)$.
\end{problem}

\begin{problem}
Prove the Generalized Distributive Law for intersection over arbitrary union: $B \cup \left(\bigcap_{i \in I} A_i\right) = \bigcap_{i \in I} (B \cup A_i)$.
\end{problem}

\begin{problem}
Let $A_n = [0, \frac{1}{n})$ for $n \in \mathbb{N}^+$. Determine the set $\bigcap_{n=1}^\infty A_n$ and rigorously prove your result using element-chasing.
\end{problem}

\subsection{Probability Spaces and Measure Theory}
\begin{problem}
Let $\mathbb{P}$ be a probability measure on a $\sigma$-algebra $\mathcal{F}$. Using the set-theoretic property that $A \cup B$ can be written as the disjoint union $A \cup (B - A)$, derive the probability measure formula: $\mathbb{P}(A \cup B) = \mathbb{P}(A) + \mathbb{P}(B) - \mathbb{P}(A \cap B)$.
\end{problem}

\begin{problem}
Using the axioms of probability and set theory, prove Boole's Inequality (the Union Bound) for two events: $\mathbb{P}(A \cup B) \le \mathbb{P}(A) + \mathbb{P}(B)$.
\end{problem}

\begin{problem}
Let $A$ and $B$ be events in a probability space. Prove that $\mathbb{P}(A \cap B) \ge \mathbb{P}(A) + \mathbb{P}(B) - 1$.
\end{problem}

\begin{problem}
Prove that if $\mathbb{P}(A) = 1$ and $\mathbb{P}(B) = 1$, then $\mathbb{P}(A \cap B) = 1$.
\end{problem}

\subsection{Quantitative Finance Modeling}
\begin{problem}
Let the state space of a market be $\Omega = \{\omega_1, \omega_2, \omega_3, \omega_4\}$. Let $A = \{\omega_1, \omega_2\}$ be the event that Tech stocks increase, and $B = \{\omega_2, \omega_3\}$ be the event that Interest Rates increase. A derivative pays $\$500$ if Tech stocks increase AND Interest Rates decrease. Find the payoff event set.
\end{problem}

\begin{problem}
Continuing from Problem 24, assume a risk-neutral probability measure $\mathbb{Q}$ where $\mathbb{Q}(\{\omega_1\}) = 0.3$, $\mathbb{Q}(\{\omega_2\}) = 0.2$, $\mathbb{Q}(\{\omega_3\}) = 0.4$, and $\mathbb{Q}(\{\omega_4\}) = 0.1$. The risk-free rate is $r = 0.04$ and time to maturity is $T = 2$. Calculate the expected discounted payoff $V_0$ of the derivative.
\end{problem}

\begin{problem}
In stochastic calculus, information is modeled as a filtration $\mathcal{F}_t$. If $\mathcal{F}_1 \subseteq \mathcal{F}_2$, prove that for any event $E \in \mathcal{F}_1$, $E \cap \Omega \in \mathcal{F}_2$.
\end{problem}

\subsection{Machine Learning Modeling}
\begin{problem}
A decision tree partitions the feature space $\mathbb{R}^2$. The first split creates set $A = \{(x,y) \in \mathbb{R}^2 : x > 5\}$. The second split, applied only to the complement of $A$, creates set $B = \{(x,y) \in \mathbb{R}^2 : y \le 10\}$. Define the set representing the leaf node for $\overline{A} \cap B$ and describe it geometrically.
\end{problem}

\begin{problem}
Prove that any single binary split of a feature space $\Omega$ into a set $C$ and its complement $\overline{C}$ creates a mutually exclusive and collectively exhaustive partition of $\Omega$.
\end{problem}

\subsection{Advanced Synthesis}
\begin{problem}
Prove that for any sets $A$ and $B$, $(A \cap B) \cup (A \cap \overline{B}) = A$. State the logical law this represents.
\end{problem}

\begin{problem}
Prove that $A \Delta B = \emptyset \iff A = B$.
\end{problem}

\newpage

% ==========================================
% SECTION 3: DETAILED SOLUTIONS
% ==========================================
\section{Detailed Solutions}
\label{sec:solutions}

\begin{proof}[Solution to Problem 1]
\textit{Proof Technique: Direct Proof via Element-Chasing.} \\
\textbf{Part 1: Show $A \cup (A \cap B) \subseteq A$.} \\
Let $x \in A \cup (A \cap B)$. By the definition of union, $x \in A$ or $x \in A \cap B$. 
Case 1: If $x \in A$, then trivially $x \in A$. 
Case 2: If $x \in A \cap B$, by the definition of intersection, $x \in A$ and $x \in B$. Thus, $x \in A$. 
In both cases, $x \in A$. Therefore, $A \cup (A \cap B) \subseteq A$. \\
\textbf{Part 2: Show $A \subseteq A \cup (A \cap B)$.} \\
Let $y \in A$. By the definition of union, for any set $S$, if $y \in A$, then $y \in A \cup S$. Let $S = (A \cap B)$. Thus, $y \in A \cup (A \cap B)$. Therefore, $A \subseteq A \cup (A \cap B)$. \\
Since both subset inclusions hold, $A \cup (A \cap B) = A$.
\end{proof}

\begin{proof}[Solution to Problem 2]
\textit{Proof Technique: Direct Proof via Element-Chasing.} \\
Assume $A \subseteq B$. We must show $A \cap B \subseteq A$ and $A \subseteq A \cap B$. \\
\textbf{Part 1: Show $A \cap B \subseteq A$.} \\
Let $x \in A \cap B$. By definition of intersection, $x \in A$ and $x \in B$. This immediately implies $x \in A$. Thus, $A \cap B \subseteq A$. (Note: This holds regardless of the assumption $A \subseteq B$). \\
\textbf{Part 2: Show $A \subseteq A \cap B$.} \\
Let $y \in A$. By our premise $A \subseteq B$, if $y \in A$, then $y \in B$. Since $y \in A$ and $y \in B$, by the definition of intersection, $y \in A \cap B$. Thus, $A \subseteq A \cap B$. \\
Since mutual inclusion holds, $A \cap B = A$.
\end{proof}

\begin{proof}[Solution to Problem 3]
\textit{Proof Technique: Biconditional Logical Equivalences.} \\
Let $x$ be an arbitrary element in the universal set $U$.
\begin{align*}
x \in (A - B) - C &\iff x \in (A - B) \land x \notin C \quad &&\text{(Definition of Set Difference)} \\
&\iff (x \in A \land x \notin B) \land x \notin C \quad &&\text{(Definition of Set Difference)} \\
&\iff x \in A \land (x \notin B \land x \notin C) \quad &&\text{(Associativity of } \land \text{)} \\
&\iff x \in A \land \neg(x \in B \lor x \in C) \quad &&\text{(De Morgan's Law of Logic)} \\
&\iff x \in A \land x \notin (B \cup C) \quad &&\text{(Definition of Union)} \\
&\iff x \in A - (B \cup C) \quad &&\text{(Definition of Set Difference)}
\end{align*}
Because the logical statements are equivalent at every step, the sets contain exactly the same elements. Thus, $(A - B) - C = A - (B \cup C)$.
\end{proof}

\begin{proof}[Solution to Problem 4]
\textit{Proof Technique: Algebraic Set Manipulation.} \\
We begin with the left-hand side and apply the definition of set difference, $X - Y = X \cap \overline{Y}$.
\begin{align*}
A - (A - B) &= A \cap \overline{(A - B)} \quad &&\text{(Definition of Set Difference)} \\
&= A \cap \overline{(A \cap \overline{B})} \quad &&\text{(Definition of Set Difference)} \\
&= A \cap (\overline{A} \cup \overline{\overline{B}}) \quad &&\text{(De Morgan's Law for Sets)} \\
&= A \cap (\overline{A} \cup B) \quad &&\text{(Double Complement Law)} \\
&= (A \cap \overline{A}) \cup (A \cap B) \quad &&\text{(Distributive Law)} \\
&= \emptyset \cup (A \cap B) \quad &&\text{(Intersection with Complement is Empty)} \\
&= A \cap B \quad &&\text{(Identity Law for Union)}
\end{align*}
Thus, $A - (A - B) = A \cap B$.
\end{proof}

\begin{proof}[Solution to Problem 5]
\textit{Proof Technique: Algebraic Set Manipulation.}
\begin{align*}
(A \cap B) - (A \cap C) &= (A \cap B) \cap \overline{(A \cap C)} \quad &&\text{(Definition of Set Difference)} \\
&= (A \cap B) \cap (\overline{A} \cup \overline{C}) \quad &&\text{(De Morgan's Law)} \\
&= ((A \cap B) \cap \overline{A}) \cup ((A \cap B) \cap \overline{C}) \quad &&\text{(Distributive Law)} \\
&= (A \cap \overline{A} \cap B) \cup (A \cap B \cap \overline{C}) \quad &&\text{(Commutativity and Associativity)} \\
&= (\emptyset \cap B) \cup (A \cap (B \cap \overline{C})) \quad &&\text{(Complement Law and Associativity)} \\
&= \emptyset \cup (A \cap (B - C)) \quad &&\text{(Domination Law and Def. of Set Difference)} \\
&= A \cap (B - C) \quad &&\text{(Identity Law)}
\end{align*}
Thus, $A \cap (B - C) = (A \cap B) - (A \cap C)$.
\end{proof}

\begin{proof}[Solution to Problem 6]
\textit{Proof Technique: Algebraic Set Manipulation.}
\begin{align*}
(A \cup B) - (A \cap B) &= (A \cup B) \cap \overline{(A \cap B)} \quad &&\text{(Definition of Set Difference)} \\
&= (A \cup B) \cap (\overline{A} \cup \overline{B}) \quad &&\text{(De Morgan's Law)} \\
&= ((A \cup B) \cap \overline{A}) \cup ((A \cup B) \cap \overline{B}) \quad &&\text{(Distributive Law)} \\
&= ((A \cap \overline{A}) \cup (B \cap \overline{A})) \cup ((A \cap \overline{B}) \cup (B \cap \overline{B})) \quad &&\text{(Distributive Law)} \\
&= (\emptyset \cup (B \cap \overline{A})) \cup ((A \cap \overline{B}) \cup \emptyset) \quad &&\text{(Complement Law)} \\
&= (B \cap \overline{A}) \cup (A \cap \overline{B}) \quad &&\text{(Identity Law)} \\
&= (B - A) \cup (A - B) \quad &&\text{(Definition of Set Difference)} \\
&= (A - B) \cup (B - A) \quad &&\text{(Commutative Law of Union)}
\end{align*}
Thus, $A \Delta B = (A \cup B) - (A \cap B)$.
\end{proof}

\begin{proof}[Solution to Problem 7]
\textit{Proof Technique: Algebraic Set Manipulation.} \\
Expand the right-hand side using the definition of symmetric difference:
\begin{align*}
(A \cap B) \Delta (A \cap C) &= ((A \cap B) - (A \cap C)) \cup ((A \cap C) - (A \cap B)) \\
&= ((A \cap B) \cap \overline{(A \cap C)}) \cup ((A \cap C) \cap \overline{(A \cap B)}) \quad &&\text{(Def. of Set Difference)} \\
&= ((A \cap B) \cap (\overline{A} \cup \overline{C})) \cup ((A \cap C) \cap (\overline{A} \cup \overline{B})) \quad &&\text{(De Morgan's Law)} \\
&= (A \cap B \cap \overline{A}) \cup (A \cap B \cap \overline{C}) \cup (A \cap C \cap \overline{A}) \cup (A \cap C \cap \overline{B}) \quad &&\text{(Distributive Law)} \\
&= \emptyset \cup (A \cap B \cap \overline{C}) \cup \emptyset \cup (A \cap C \cap \overline{B}) \quad &&\text{($A \cap \overline{A} = \emptyset$)} \\
&= A \cap ((B \cap \overline{C}) \cup (C \cap \overline{B})) \quad &&\text{(Factoring out } A \cap \text{ via Distributive Law)} \\
&= A \cap ((B - C) \cup (C - B)) \quad &&\text{(Def. of Set Difference)} \\
&= A \cap (B \Delta C) \quad &&\text{(Def. of Symmetric Difference)}
\end{align*}
\end{proof}

\begin{proof}[Solution to Problem 8]
\textit{Proof Technique: Element-Chasing (Biconditional Implication).} \\
We must prove two directions. \\
\textbf{Direction 1 ($\implies$):} Assume $A \subseteq B$. We must show $\overline{B} \subseteq \overline{A}$. \\
Let $x \in \overline{B}$. By definition, $x \notin B$. Since $A \subseteq B$, if $x$ were in $A$, it would have to be in $B$. But $x \notin B$, so $x$ cannot be in $A$. Thus, $x \notin A$, which means $x \in \overline{A}$. Therefore, $\overline{B} \subseteq \overline{A}$. \\
\textbf{Direction 2 ($\impliedby$):} Assume $\overline{B} \subseteq \overline{A}$. We must show $A \subseteq B$. \\
Let $y \in A$. By definition, $y \notin \overline{A}$. Since $\overline{B} \subseteq \overline{A}$, if $y$ were in $\overline{B}$, it would have to be in $\overline{A}$. But $y \notin \overline{A}$, so $y$ cannot be in $\overline{B}$. Thus, $y \notin \overline{B}$, which means $y \in B$. Therefore, $A \subseteq B$. \\
Thus, $A \subseteq B \iff \overline{B} \subseteq \overline{A}$.
\end{proof}

\begin{proof}[Solution to Problem 9]
\textit{Proof Technique: Element-Chasing.} \\
The logical tautology $P \implies (P \lor Q)$ translates to the set-theoretic statement: $A \subseteq (A \cup B)$. \\
Let $x \in A$. By the definition of union, $A \cup B = \{y : y \in A \lor y \in B\}$. Since $x \in A$ is true, the disjunction $x \in A \lor x \in B$ is logically true regardless of whether $x \in B$. Therefore, $x \in A \cup B$. Since $x$ was arbitrary, $A \subseteq (A \cup B)$.
\end{proof}

\begin{proof}[Solution to Problem 10]
\textit{Proof Technique: Biconditional Logical Equivalences.} \\
Let $x \in U$.
\begin{align*}
x \in \overline{A \cup (B \cap C)} &\iff \neg(x \in A \cup (B \cap C)) \quad &&\text{(Definition of Complement)} \\
&\iff \neg(x \in A \lor (x \in B \land x \in C)) \quad &&\text{(Definition of Union and Intersection)} \\
&\iff x \notin A \land \neg(x \in B \land x \in C) \quad &&\text{(De Morgan's Law of Logic)} \\
&\iff x \notin A \land (x \notin B \lor x \notin C) \quad &&\text{(De Morgan's Law of Logic)} \\
&\iff (x \notin A \land x \notin B) \lor (x \notin A \land x \notin C) \quad &&\text{(Distributive Law of Logic)} \\
&\iff (x \in \overline{A} \land x \in \overline{B}) \lor (x \in \overline{A} \land x \in \overline{C}) \quad &&\text{(Definition of Complement)} \\
&\iff (x \in \overline{A} \cap \overline{B}) \lor (x \in \overline{A} \cap \overline{C}) \quad &&\text{(Definition of Intersection)} \\
&\iff x \in (\overline{A} \cap \overline{B}) \cup (\overline{A} \cap \overline{C}) \quad &&\text{(Definition of Union)}
\end{align*}
Thus, the sets are equal.
\end{proof}

\begin{proof}[Solution to Problem 11]
\textit{Proof Technique: Counterexample.} \\
Let $A = \{1\}$, $B = \{1\}$, and $C = \{1\}$. \\
Evaluate the left-hand side (LHS): $A - (B - C) = \{1\} - (\{1\} - \{1\}) = \{1\} - \emptyset = \{1\}$. \\
Evaluate the right-hand side (RHS): $(A - B) - C = (\{1\} - \{1\}) - \{1\} = \emptyset - \{1\} = \emptyset$. \\
Since $\{1\} \neq \emptyset$, $A - (B - C) \neq (A - B) - C$. Set difference is not associative.
\end{proof}

\begin{proof}[Solution to Problem 12]
\textit{Proof Technique: Counterexample.} \\
Let $A = \{1\}$, $B = \emptyset$, and $C = \{1\}$. \\
Evaluate the LHS: $(A \cup B) - C = (\{1\} \cup \emptyset) - \{1\} = \{1\} - \{1\} = \emptyset$. \\
Evaluate the RHS: $A \cup (B - C) = \{1\} \cup (\emptyset - \{1\}) = \{1\} \cup \emptyset = \{1\}$. \\
Since $\emptyset \neq \{1\}$, the claim is false.
\end{proof}

\begin{proof}[Solution to Problem 13]
\textit{Proof Technique: Logical Analysis.} \\
\textbf{Flaw:} The statement commits an Improper Negation fallacy. The condition $x \notin A \cap B$ means $\neg(x \in A \land x \in B)$. By De Morgan's Law for logic, $\neg(P \land Q) \iff \neg P \lor \neg Q$. The statement incorrectly distributed the negation as an AND ($\land$) instead of an OR ($\lor$). \\
\textbf{Correct Negation:} ``If $x \notin A \cap B$, then $x \notin A \lor x \notin B$.''
\end{proof}

\begin{proof}[Solution to Problem 14]
\textit{Proof Technique: Logical Analysis and Counterexample.} \\
\textbf{Flaw:} The proof assumes that if every element of $A$ belongs to $B \cup C$, then all elements of $A$ must belong entirely to $B$ or entirely to $C$. This is a quantifier scope fallacy: $\forall x \in A, (x \in B \lor x \in C)$ does not imply $(\forall x \in A, x \in B) \lor (\forall x \in A, x \in C)$. \\
\textbf{Counterexample:} Let $A = \{1, 2\}$, $B = \{1\}$, $C = \{2\}$. $B \cup C = \{1, 2\}$. Here, $A \subseteq B \cup C$ is true ($\{1, 2\} \subseteq \{1, 2\}$). However, $A \subseteq B$ is false ($\{1, 2\} \not\subseteq \{1\}$) and $A \subseteq C$ is false ($\{1, 2\} \not\subseteq \{2\}$).
\end{proof}

\begin{proof}[Solution to Problem 15]
\textit{Proof Technique: Counterexample.} \\
The claim is false. Let $A = \{1\}$ and $B = \{2\}$. \\
Then $A \cup B = \{1, 2\}$. \\
The power set of the union is $\mathcal{P}(A \cup B) = \{\emptyset, \{1\}, \{2\}, \{1, 2\}\}$. \\
The individual power sets are $\mathcal{P}(A) = \{\emptyset, \{1\}\}$ and $\mathcal{P}(B) = \{\emptyset, \{2\}\}$. \\
Their union is $\mathcal{P}(A) \cup \mathcal{P}(B) = \{\emptyset, \{1\}, \{2\}\}$. \\
Notice that $\{1, 2\} \in \mathcal{P}(A \cup B)$, but $\{1, 2\} \notin \mathcal{P}(A) \cup \mathcal{P}(B)$. Thus, the sets are not equal.
\end{proof}

\begin{proof}[Solution to Problem 16]
\textit{Proof Technique: Biconditional Logical Equivalences for Arbitrary Families.} \\
Let $x \in U$.
\begin{align*}
x \in \overline{\bigcap_{i \in I} A_i} &\iff x \notin \bigcap_{i \in I} A_i \quad &&\text{(Definition of Complement)} \\
&\iff \neg\left(\forall i \in I, x \in A_i\right) \quad &&\text{(Definition of Arbitrary Intersection)} \\
&\iff \exists i \in I, x \notin A_i \quad &&\text{(Negation of Universal Quantifier)} \\
&\iff \exists i \in I, x \in \overline{A_i} \quad &&\text{(Definition of Complement)} \\
&\iff x \in \bigcup_{i \in I} \overline{A_i} \quad &&\text{(Definition of Arbitrary Union)}
\end{align*}
Thus, $\overline{\bigcap_{i \in I} A_i} = \bigcup_{i \in I} \overline{A_i}$.
\end{proof}

\begin{proof}[Solution to Problem 17]
\textit{Proof Technique: Biconditional Logical Equivalences.} \\
Let $x \in U$.
\begin{align*}
x \in B \cap \left(\bigcup_{i \in I} A_i\right) &\iff x \in B \land x \in \bigcup_{i \in I} A_i \quad &&\text{(Definition of Intersection)} \\
&\iff x \in B \land (\exists i \in I, x \in A_i) \quad &&\text{(Definition of Arbitrary Union)} \\
&\iff \exists i \in I, (x \in B \land x \in A_i) \quad &&\text{(Distributing independent statement into } \exists \text{)} \\
&\iff \exists i \in I, x \in (B \cap A_i) \quad &&\text{(Definition of Intersection)} \\
&\iff x \in \bigcup_{i \in I} (B \cap A_i) \quad &&\text{(Definition of Arbitrary Union)}
\end{align*}
Thus, the sets are equal.
\end{proof}

\begin{proof}[Solution to Problem 18]
\textit{Proof Technique: Biconditional Logical Equivalences.} \\
Let $x \in U$.
\begin{align*}
x \in B \cup \left(\bigcap_{i \in I} A_i\right) &\iff x \in B \lor x \in \bigcap_{i \in I} A_i \quad &&\text{(Definition of Union)} \\
&\iff x \in B \lor (\forall i \in I, x \in A_i) \quad &&\text{(Definition of Arbitrary Intersection)} \\
&\iff \forall i \in I, (x \in B \lor x \in A_i) \quad &&\text{(Distributing independent statement into } \forall \text{)} \\
&\iff \forall i \in I, x \in (B \cup A_i) \quad &&\text{(Definition of Union)} \\
&\iff x \in \bigcap_{i \in I} (B \cup A_i) \quad &&\text{(Definition of Arbitrary Intersection)}
\end{align*}
Thus, the sets are equal.
\end{proof}

\begin{proof}[Solution to Problem 19]
\textit{Proof Technique: Element-Chasing.} \\
\textbf{Claim:} $\bigcap_{n=1}^\infty [0, \frac{1}{n}) = \{0\}$. \\
\textbf{Part 1: Show $\{0\} \subseteq \bigcap_{n=1}^\infty A_n$.} \\
Let $x = 0$. For every $n \in \mathbb{N}^+$, $0 \le 0 < \frac{1}{n}$. Thus, $0 \in [0, \frac{1}{n})$ for all $n$. By definition of arbitrary intersection, $0 \in \bigcap_{n=1}^\infty A_n$. \\
\textbf{Part 2: Show $\bigcap_{n=1}^\infty A_n \subseteq \{0\}$.} \\
Let $x \in \bigcap_{n=1}^\infty A_n$. This means for all $n \in \mathbb{N}^+$, $0 \le x < \frac{1}{n}$. 
Assume for contradiction that $x > 0$. By the Archimedean property of real numbers, there exists an integer $N > 0$ such that $\frac{1}{N} < x$. But this contradicts the premise that $x < \frac{1}{n}$ for all $n$. Thus, $x$ cannot be strictly greater than 0. Since $x \ge 0$, we must have $x = 0$. Thus, $x \in \{0\}$. \\
Therefore, the intersection is exactly $\{0\}$.
\end{proof}

\begin{proof}[Solution to Problem 20]
\textit{Proof Technique: Measure Theory Axioms.} \\
By set theory, $A \cup B = A \cup (B - A)$. Furthermore, $A$ and $B - A$ are disjoint because $A \cap (B - A) = A \cap (B \cap \overline{A}) = (A \cap \overline{A}) \cap B = \emptyset \cap B = \emptyset$. \\
By the additivity axiom of probability for disjoint sets:
\[ \mathbb{P}(A \cup B) = \mathbb{P}(A) + \mathbb{P}(B - A) \]
We also know that $B = (A \cap B) \cup (B - A)$. Since $(A \cap B)$ and $(B - A)$ are disjoint (as $A \cap \overline{A} = \emptyset$), we apply the additivity axiom again:
\[ \mathbb{P}(B) = \mathbb{P}(A \cap B) + \mathbb{P}(B - A) \]
Rearranging this gives $\mathbb{P}(B - A) = \mathbb{P}(B) - \mathbb{P}(A \cap B)$. \\
Substitute this back into the first equation:
\[ \mathbb{P}(A \cup B) = \mathbb{P}(A) + \mathbb{P}(B) - \mathbb{P}(A \cap B) \]
\end{proof}

\begin{proof}[Solution to Problem 21]
\textit{Proof Technique: Algebraic Manipulation of Probability Measure.} \\
From Problem 20, we established:
\[ \mathbb{P}(A \cup B) = \mathbb{P}(A) + \mathbb{P}(B) - \mathbb{P}(A \cap B) \]
By Kolmogorov's First Axiom, the probability of any event is non-negative. Therefore, $\mathbb{P}(A \cap B) \ge 0$. \\
Subtracting a non-negative number from $\mathbb{P}(A) + \mathbb{P}(B)$ means:
\[ \mathbb{P}(A \cup B) \le \mathbb{P}(A) + \mathbb{P}(B) \]
This completes the proof of Boole's Inequality for two sets.
\end{proof}

\begin{proof}[Solution to Problem 22]
\textit{Proof Technique: Algebraic Manipulation of Probability Measure.} \\
From Problem 20, we have:
\[ \mathbb{P}(A \cup B) = \mathbb{P}(A) + \mathbb{P}(B) - \mathbb{P}(A \cap B) \]
Rearranging to solve for $\mathbb{P}(A \cap B)$ yields:
\[ \mathbb{P}(A \cap B) = \mathbb{P}(A) + \mathbb{P}(B) - \mathbb{P}(A \cup B) \]
By Kolmogorov's Axioms, the probability of any event is at most 1. Thus, $\mathbb{P}(A \cup B) \le 1$. \\
Subtracting a number less than or equal to 1 means we are subtracting at most 1, which gives the inequality:
\[ \mathbb{P}(A \cap B) \ge \mathbb{P}(A) + \mathbb{P}(B) - 1 \]
\end{proof}

\begin{proof}[Solution to Problem 23]
\textit{Proof Technique: Direct Application of Probability Bounds.} \\
From Problem 22, we established $\mathbb{P}(A \cap B) \ge \mathbb{P}(A) + \mathbb{P}(B) - 1$. \\
Substitute $\mathbb{P}(A) = 1$ and $\mathbb{P}(B) = 1$:
\[ \mathbb{P}(A \cap B) \ge 1 + 1 - 1 = 1 \]
Since probability cannot exceed 1, $\mathbb{P}(A \cap B) \le 1$. \\
Therefore, $\mathbb{P}(A \cap B) = 1$.
\end{proof}

\begin{proof}[Solution to Problem 24]
\textit{Proof Technique: Set Operations in Concrete Scenarios.} \\
The event that Tech stocks increase is $A = \{\omega_1, \omega_2\}$. \\
The event that Interest Rates increase is $B = \{\omega_2, \omega_3\}$. \\
The event that Interest Rates decrease (assuming binary states for simplicity, or simply the complement in this $\sigma$-algebra) is $\overline{B} = \Omega - B = \{\omega_1, \omega_4\}$. \\
The derivative pays off if Tech increases AND Rates decrease. This is the intersection:
\[ A \cap \overline{B} = \{\omega_1, \omega_2\} \cap \{\omega_1, \omega_4\} = \{\omega_1\} \]
The payoff event set is $\{\omega_1\}$.
\end{proof}

\begin{proof}[Solution to Problem 25]
\textit{Proof Technique: Quantitative Finance Calculation.} \\
The payoff event is $E = \{\omega_1\}$. \\
The risk-neutral probability of this event is $\mathbb{Q}(E) = \mathbb{Q}(\{\omega_1\}) = 0.3$. \\
The expected payoff under $\mathbb{Q}$ is $500 \times \mathbb{Q}(E) = 500 \times 0.3 = 150$. \\
The expected discounted payoff $V_0$ is calculated using continuous compounding:
\[ V_0 = e^{-rT} \mathbb{E}^\mathbb{Q}[\text{Payoff}] = e^{-0.04 \times 2} \times 150 = 150 e^{-0.08} \]
\end{proof}

\begin{proof}[Solution to Problem 26]
\textit{Proof Technique: Direct Proof via Set Identities.} \\
Assume $\mathcal{F}_1 \subseteq \mathcal{F}_2$ and $E \in \mathcal{F}_1$. \\
By the definition of a subset, if $E \in \mathcal{F}_1$, then $E \in \mathcal{F}_2$. \\
By the Identity Law of set intersection, for any set $E$ within a universal set $\Omega$, $E \cap \Omega = E$. \\
Since $E \in \mathcal{F}_2$ and $E \cap \Omega = E$, it follows by substitution that $E \cap \Omega \in \mathcal{F}_2$.
\end{proof}

\begin{proof}[Solution to Problem 27]
\textit{Proof Technique: Set-Builder Notation and Geometric Interpretation.} \\
The first set is $A = \{(x,y) : x > 5\}$. Its complement is $\overline{A} = \{(x,y) : x \le 5\}$. \\
The second set is $B = \{(x,y) : y \le 10\}$. \\
The leaf node represents the intersection $\overline{A} \cap B$:
\[ \overline{A} \cap B = \{(x,y) \in \mathbb{R}^2 : x \le 5 \land y \le 10\} \]
Geometrically, this represents the bottom-left quadrant of the Cartesian plane originating from the point $(5, 10)$, including the boundary lines $x=5$ and $y=10$.
\end{proof}

\begin{proof}[Solution to Problem 28]
\textit{Proof Technique: Direct Proof of Partition Properties.} \\
A partition of a space $\Omega$ requires the subsets to be mutually exclusive (disjoint) and collectively exhaustive (their union is $\Omega$). \\
Let the split create sets $C$ and $\overline{C}$. \\
\textbf{Mutually Exclusive:} By the definition of complement, $\overline{C} = \{x \in \Omega : x \notin C\}$. Thus, there is no element $x$ such that $x \in C$ and $x \in \overline{C}$. Therefore, $C \cap \overline{C} = \emptyset$. \\
\textbf{Collectively Exhaustive:} By the Law of Excluded Middle in logic, for any $x \in \Omega$, either $x \in C$ or $x \notin C$. Thus, $x \in C \cup \overline{C}$. Therefore, $C \cup \overline{C} = \Omega$. \\
Since both conditions are met, the split creates a valid partition.
\end{proof}

\begin{proof}[Solution to Problem 29]
\textit{Proof Technique: Algebraic Set Manipulation.}
\begin{align*}
(A \cap B) \cup (A \cap \overline{B}) &= A \cap (B \cup \overline{B}) \quad &&\text{(Distributive Law)} \\
&= A \cap U \quad &&\text{(Complement Law: } B \cup \overline{B} = U \text{)} \\
&= A \quad &&\text{(Identity Law)}
\end{align*}
This represents the \textbf{Law of Total Probability} in measure theory (or the Resolution rule in logic), where a set can be reconstructed by partitioning it against any other set $B$.
\end{proof}

\begin{proof}[Solution to Problem 30]
\textit{Proof Technique: Biconditional Logical Equivalences.} \\
We must prove $A \Delta B = \emptyset \iff A = B$. \\
\textbf{Direction 1 ($\implies$):} Assume $A \Delta B = \emptyset$. \\
By definition, $(A - B) \cup (B - A) = \emptyset$. \\
The union of two sets is empty if and only if both sets are empty. Thus, $A - B = \emptyset$ and $B - A = \emptyset$. \\
$A - B = \emptyset \implies A \subseteq B$. \\
$B - A = \emptyset \implies B \subseteq A$. \\
By the Axiom of Extensionality, $A \subseteq B$ and $B \subseteq A \implies A = B$. \\
\textbf{Direction 2 ($\impliedby$):} Assume $A = B$. \\
Then $A - B = A - A = A \cap \overline{A} = \emptyset$. \\
Similarly, $B - A = B - B = \emptyset$. \\
Thus, $A \Delta B = \emptyset \cup \emptyset = \emptyset$. \\
Therefore, $A \Delta B = \emptyset \iff A = B$.
\end{proof}

\newpage

% ==========================================
% SECTION 4: CONCEPT CHECKLIST (REFERENCE)
% ==========================================
\section{Concept Checklist \& Question Mapping (Reference)}
\label{sec:checklist_end}

\begin{itemize}
    \item \textbf{C1: Element-Chasing Proofs (Mutual Inclusion):} Rigorous deduction of $X \subseteq Y$ and $Y \subseteq X$. (Problems 1, 2, 8, 9)
    \item \textbf{C2: Biconditional Logical Equivalences:} Chaining $\iff$ statements using Boolean logic. (Problems 3, 10)
    \item \textbf{C3: Algebraic Set Manipulation:} Utilizing established axioms (De Morgan's, Distributive) as algebraic variables. (Problems 4, 5, 6, 7)
    \item \textbf{C4: Set Difference Translations:} Converting $A - B$ to $A \cap \overline{B}$. (Problems 4, 5, 6, 11, 12)
    \item \textbf{C5: Propositional Logic to Set Theory Isomorphism:} Translating tautologies and logical laws to set identities. (Problems 13, 14, 15)
    \item \textbf{C6: Infinite Families of Sets:} Arbitrary unions ($\bigcup_{i \in I}$) and intersections ($\bigcap_{i \in I}$). (Problems 16, 17, 18, 19)
    \item \textbf{C7: Generalized De Morgan's Laws:} Extending finite complements to arbitrary index sets. (Problem 16)
    \item \textbf{C8: Probability Spaces \& Measure Theory ($\sigma$-algebras):} Set-theoretic foundations for mutually exclusive events and probability bounds. (Problems 20, 21, 22, 23)
    \item \textbf{C9: Quantitative Finance Modeling:} Binomial pricing, risk-neutral measures ($\mathbb{Q}$), and expected discounted payoffs. (Problems 24, 25, 26)
    \item \textbf{C10: Machine Learning Modeling:} Feature space partitioning and decision tree logic. (Problems 27, 28)
    \item \textbf{C11: Diagnostic Flaws \& Counterexamples:} Identifying inverse fallacies, non-associativity, and vacuous truth errors. (Problems 29, 30)
\end{itemize}

\end{document}